% ============================================================
% COLLISION AVOIDANCE  (Requirement 4c)
% ============================================================
\section{Collision Avoidance}
\label{sec:avoidance}

\subsection{Working Principle}
Obstacles are avoided in two layers (Fig.~\ref{fig:layers}). The first
is deliberative: when an obstacle is detected near the path still to be
driven, that part of the path is planned again around it. The second is
reactive: a potential field \cite{khatib1986} steers the vehicle away
from any obstacle whose centre comes within $0.8$~m. Re-planning does most of
the work because it acts early and keeps the guarantees of the planner;
the reactive layer covers what re-planning cannot, such as an obstacle
confirmed too late. Neither layer reads the obstacle list. Every
obstacle position they use comes from the onboard detector.

\begin{figure}[!t]
    \centering
    \begin{tikzpicture}[
        font=\footnotesize,
        node distance=0.38cm,
        io/.style={align=center, inner sep=1pt},
        st/.style={draw, rounded corners=2pt, minimum height=0.6cm,
                   text width=3.4cm, align=center, fill=black!4},
        arr/.style={-{Stealth[length=1.8mm]}, semithick}]
        % layer 1
        \node[io, text width=3.5cm] (i1) {detections, $\hat{\mathbf{x}}$,\\waypoint index};
        \node[st, below=of i1] (m)  {Remember and confirm obstacles};
        \node[st, below=of m]  (t)  {Leg within $1.0$~m\\of an obstacle?};
        \node[st, below=of t]  (r)  {A* and PSO with\\obstacle keep-outs};
        \node[io, below=of r, text width=3.5cm] (o1) {new path,\\to guidance};
        \draw[arr] (i1) -- (m);  \draw[arr] (m) -- (t);
        \draw[arr] (t) -- (r);   \draw[arr] (r) -- (o1);
        % layer 2
        \node[io, text width=3.5cm, right=0.5cm of i1] (i2) {$\theta_{\text{cmd}}$, detections,\\lidar ranges, $d_{\text{wp}}$};
        \node[st, below=of i2] (p)  {Attraction plus obstacle and wall repulsion};
        \node[st, below=of p]  (f)  {Low-pass filter\\on heading};
        \node[st, below=of f]  (s)  {Speed factor};
        \node[io, below=of s, text width=3.5cm] (o2) {$\theta_{\text{safe}}$, $v_{\text{scale}}$,\\to controller};
        \draw[arr] (i2) -- (p);  \draw[arr] (p) -- (f);
        \draw[arr] (f) -- (s);   \draw[arr] (s) -- (o2);
        \node[above=0.1cm of i1, font=\footnotesize\bfseries] {Layer 1: re-planning};
        \node[above=0.1cm of i2, font=\footnotesize\bfseries] {Layer 2: reactive};
    \end{tikzpicture}
    \caption{The two avoidance layers: inputs, stages and outputs.}
    \label{fig:layers}
\end{figure}

\textbf{Layer 1: online re-planning.} The re-planner runs at $100$~Hz.
\begin{enumerate}
    \item \textit{Remember.} Each detection is converted to a world
    position with Eq.~\eqref{eq:detection}. A detection within $0.75$~m
    of a known obstacle updates that obstacle's position with a
    first-order filter; any other starts a new obstacle.
    \item \textit{Confirm.} An obstacle is trusted once it has been
    seen on $20$ re-planner steps ($0.2$~s, two or three detector updates).
    \item \textit{Trigger.} A leg is re-planned if the part still to be
    driven passes within $1.0$~m of a confirmed obstacle's centre. Re-planning
    runs at most once per second, and only when an obstacle is new or
    its position has moved by more than $0.3$~m.
    \item \textit{Re-plan.} Each confirmed obstacle is added to the
    planner's clearance field as a keep-out disc of radius $0.8$~m about
    its centre, which the planner treats like an inflated wall. A*
    routes the leg again, from the vehicle's position for the leg in
    progress, and PSO smooths it. A* first tries a $1.0$~m keep-out, so
    that it prefers the open side of an obstacle to a gap between the
    obstacle and a wall.
    \item \textit{Hand over.} The new path is published to guidance.
\end{enumerate}
The $0.8$~m keep-out, measured from the obstacle's centre, is the
$0.55$~m that the $0.40$~m clearance requires there, plus $0.25$~m for
detection error and corner cutting.
Figure~\ref{fig:replan} shows a re-plan.

\begin{figure}[!t]
    \centering
    \reportfig{fig_replan_k3.png}
    \caption{A re-plan in mission 22. The path before the re-plan
    (dashed) ran through an obstacle missing from the map. The diamond
    marks the vehicle's position when the obstacle was confirmed,
    $4.8$~m before it; the vehicle then follows the re-planned path
    outside the keep-out.}
    \label{fig:replan}
\end{figure}

\textbf{Layer 2: reactive avoidance.} The guidance heading is treated
as a unit attraction vector in the body frame. Each remembered obstacle
at range $\rho_{o}$ and bearing $\phi_{o}$ from the vehicle, closer than
the reach $d_{r} = 0.8$~m, adds a repulsion of magnitude
\begin{equation}
    w_{o} = k_{\text{rep}}\left(\frac{1}{\rho_{o}} - \frac{1}{d_{r}}\right)
    \label{eq:repulsion}
\end{equation}
directed away from it, and a tangential term of magnitude
$w_{o}k_{\text{tan}}\cos\phi_{o}$ that slides the vehicle around the
obstacle on the side the guidance heading points to. The side is chosen
once per encounter and held, which prevents oscillation in front of an
obstacle. Obstacles more than $108^{\circ}$ off the heading have been
passed and are ignored.
Each lidar beam shorter than $1.2$~m adds a repulsion of the same form,
with gain $k_{\text{wall}}$, away from the wall it hit. The gains are
$k_{\text{rep}} = k_{\text{tan}} = 1.0$ and $k_{\text{wall}} = 1.8$. The
heading of the summed vector is low-pass filtered with a time constant
of about $0.5$~s and sent to the controller as $\theta_{\text{safe}}$,
with a speed factor $v_{\text{scale}} \in [0.3, 1]$ that slows the
vehicle near an obstacle and in sharp turns. The lower limit is not
zero because a steered vehicle must keep moving to turn.

\subsection{Sensor Configuration and Obstacle Detection}
Two sensors are used, both at the vehicle's centre with a range of
$5$~m (Fig.~\ref{fig:sensing}). The lidar has $13$ beams over
$180^{\circ}$, updates at $20$~Hz, sees walls only, and is used for the
wall repulsion. The object detector has a $120^{\circ}$ field of view,
updates at $10$~Hz, and returns the range $\rho$, bearing $\phi$ and
label of up to three obstacles. A detection is placed in the world
frame from the state estimate:
\begin{equation}
\begin{aligned}
    x_{o} &= \hat{x} + \rho\cos(\hat{\theta} + \phi), \\
    y_{o} &= \hat{y} + \rho\sin(\hat{\theta} + \phi).
\end{aligned}
\label{eq:detection}
\end{equation}
Any error in the estimate therefore becomes an error in the obstacle's
position. Filtering repeated detections reduces that error, and the
margin in the keep-out absorbs what remains.

\begin{figure}[!t]
    \centering
    \begin{tikzpicture}[scale=0.56, font=\scriptsize]
        % lidar half-disc (180 deg, 5 m) and its 13 beams
        \fill[black!6] (0,0) -- (0,-5) arc[start angle=-90, end angle=90, radius=5] -- cycle;
        \foreach \a in {-90,-75,...,90}
            \draw[black!35, very thin] (0,0) -- (\a:5);
        % detector cone (120 deg, 5 m)
        \fill[CoverGold!35, opacity=0.7] (0,0) -- (-60:5) arc[start angle=-60, end angle=60, radius=5] -- cycle;
        \draw[CoverGold!80!black] (0,0) -- (-60:5) arc[start angle=-60, end angle=60, radius=5] -- cycle;
        % heading, obstacle, range and bearing
        \draw[-{Stealth[length=1.6mm]}, semithick] (0,0) -- (0:5.9) node[right] {heading};
        \fill[red!70!black] (35:3.6) circle[radius=0.15];
        \draw[red!70!black] (35:3.6) circle[radius=0.55];
        \draw[violet, dashed] (35:3.6) circle[radius=0.8];
        \draw[semithick] (0,0) -- node[above, sloped] {$\rho$} (35:3.6);
        \draw (0:1.5) arc[start angle=0, end angle=35, radius=1.5];
        \node at (17:1.85) {$\phi$};
        % vehicle
        \fill[white] (0,0) circle[radius=0.25];  \draw[semithick] (0,0) circle[radius=0.25];
        \node[align=left, anchor=west] at (0.4,-4.4) {lidar: $180^{\circ}$, 13 beams};
        \node[align=left, anchor=west] at (2.7,-2.9) {detector: $120^{\circ}$};
    \end{tikzpicture}
    \caption{Sensing geometry to scale. An obstacle is shown with its
    clearance ring (solid, $0.40$~m from its surface) and its keep-out
    (dashed, $0.8$~m from its centre).}
    \label{fig:sensing}
\end{figure}

\subsection{Integration with the Guidance Subsystem}
The two layers meet guidance at different points
(Fig.~\ref{fig:architecture}). The re-planner acts on the \emph{input} of guidance:
it replaces the path that guidance follows, and guidance continues with
the same logic on the new path. No mode switch is needed, and the
waypoints are never changed. The reactive layer acts on the
\emph{output} of guidance: it sits between guidance and the controller and modifies
the heading command and the speed.

The reactive layer also hands control back without a switch. When no
obstacle is within reach, its obstacle repulsion is zero and its output
is the guidance heading, filtered and adjusted only by any wall
repulsion. Near a waypoint the reach shrinks to the distance to the
waypoint, though never below $0.7$~m, the lidar wall reach likewise
shrinks from $1.2$~m to no less than $0.5$~m, and the tangential term
fades in proportion, so an obstacle beside a waypoint cannot hold the vehicle
away from it.

\subsection{Limitations}
\textit{Field of view.} An obstacle is known only once it has been
inside the $120^{\circ}$, $5$~m cone, so one just around a corner is
detected late and leaves the reactive layer little room.
\textit{Static obstacles.} A remembered obstacle is never forgotten and
its position is averaged, so obstacles are assumed not to move.
\textit{Estimate error.} Obstacle positions inherit the error of the
state estimate. The margins are sized against its measured
distribution (Section~\ref{sec:estimator}) but are fixed; they do not
shrink or grow with the filter's covariance during a mission. \textit{Local minima.} The reactive layer can stall where
repulsion balances attraction, for example in a gap narrower than the
keep-out; re-planning removes most such cases by routing round the open
side. \textit{Computation.} A re-plan took $0.22$~s to compute on
average and up to $0.92$~s over the $25$ missions of
Section~\ref{sec:results}. The simulation clock waits for it, so it is
not counted in the mission time; a real vehicle would travel up to
about $0.9$~m on the old path first.
